\label{res:rank}
The retained condition is $k\nmid r$. At $k=6$, for example, the residue
$r=2$ is retained. At base $2$ and level $3$, the fifteen sections have nine
coordinates; the omitted positive-residue sections satisfy
\[
 \ph(8n+2)=\ph(4n+1),\qquad
 \ph(8n+4)=2\ph(2n+1),\qquad
 \ph(8n+6)=\ph(4n+3).
\]

\subsection{The evaluation matrix}\label{sec:rank}
The two ingredients are scalar reduction and arithmetic independence.
The local formula for $\ph$ gives the displayed reductions. For a prime
$p\mid k$ that does not divide $u$, the argument $k^{j-t}n+u$ is a $p$-adic
unit, so multiplication by $k^t$ contributes its power of $p$ and one factor
$1-1/p$. When $p\mid u$, the argument already contains $p$, and only its
power contributes. This gives $C_k(t,u)$.

For independence, consider finitely many forms $L_i(n)=a_in+b_i$ with
positive slopes and $a_ib_j\ne a_jb_i$ for $i\ne j$. A finite shift makes
the intercepts positive. Write $L_i=g_iA_i$, where $A_i$ is primitive.
Choose an odd prime $\ell$ avoiding the slopes, all $g_i\ph(g_i)$, and the
nonzero cross determinants. In row $i$, choose distinct primes
$q_{ij}\equiv1\pmod\ell$ for $j\ne i$, outside the same finite exceptional
set, and impose
\[
 L_j(n)\equiv0\pmod{q_{ij}},\qquad A_i(n)\equiv2\pmod\ell.
\]
The Chinese remainder theorem combines these congruences. The resulting
progression for $A_i(n)$ is reduced: a prime $q_{ij}$ dividing $A_i(n)$
would also divide the nonzero cross determinant; $\ell$ is excluded by the
chosen residue; and primitivity excludes the primes of the slope.
Dirichlet's theorem therefore supplies arbitrarily large primes $p=A_i(n)$
with $p\nmid g_i$. The diagonal value is
$\ph(g_i)(p-1)\not\equiv0\pmod\ell$. Each off-diagonal value is divisible
by $\ell$, since $q_{ij}-1$ divides the corresponding totient.
One evaluation per row gives
\[
 \bigl(\ph(L_j(n_i))\bigr)_{i,j}
 \equiv\operatorname{diag}(u_1,\ldots,u_s)\pmod\ell,
 \qquad u_i\ne0.
\]
The determinant is nonzero over $\Q$.

Apply this construction after restricting a relation among the retained
sections to $n=km+1$. The positive-residue sections give pairwise
nonproportional forms
\[
 k^{j+1}m+(k^j+r),\qquad k\nmid r,
\]
and the remaining form is $km+1$. Equality between two affine fractions
would make the residue at the larger level divisible by $k$. The matrix
therefore eliminates every positive-residue coefficient. If $A,B$ are the
two zero-residue coefficients, the restriction gives
$A+B\ph(k)=0$. Evaluation at $n=k$ gives $A+Bk=0$, since
$\ph(k^2)=k\ph(k)$. As $\ph(k)<k$, both coefficients vanish.
The reductions give spanning and the count is
\[
 2+\sum_{j=1}^e(k^j-k^{j-1})=k^e+1.
\]
This proves Theorem~\ref{thm:kkernelrank}.

\begin{corollary}[Dyadic basis]\label{res:basis}
The family
\[
 \{\ph_{0,0},\ph_{1,0}\}
 \cup\{\ph_{j,r}:j\ge1,\ 0<r<2^j,\ r\text{ odd}\},
 \qquad \ph_{j,r}(n)=\ph(2^jn+r),
\]
is a basis for the rational span of the full dyadic kernel. Every rational
relation is generated by the scalar reductions. The complete level-zero
truncation has dimension one.
\end{corollary}
\begin{proof}
Every finite subfamily lies in a sufficiently high truncation. The theorem
gives its independence; the reductions give equality of the full spans.
\end{proof}

\begin{corollary}[Integral normal form]\label{cor:integral-normal-form}
The canonical family is also a $\Z$-basis of the module generated by the
sections through level $e$. Its integral relation module has a basis with
one elementary reduction per omitted channel and rank
$\sum_{j=1}^{e-1}k^j$.
\end{corollary}
\begin{proof}
The denominator $\prod_{p\mid k,\,p\nmid u}p$ in $C_k(t,u)$ divides $k$,
so all reduction scalars are integers. In the free module on the channels,
subtracting one elementary reduction for each omitted coordinate leaves a
relation among the independent retained coordinates. The omitted-coordinate
pivots are all one, which also gives independence of the elementary relations.
\end{proof}
The module here is the one generated by the channels. It is not saturated
in the group of all integer-valued sequences: at base $2$ and $e\ge2$,
$\ph(4n+3)/2$ is integer-valued but has canonical coefficient $1/2$.

\subsection{A reusable consequence}
\begin{corollary}[Freezing periodic coefficients]\label{cor:periodic-freezing}
If $L_i$ are pairwise nonproportional positive-slope affine forms and $w_i$
are rational-valued periodic sequences, then
\[
 \sum_i w_i(n)\ph(L_i(n))=0\quad\text{for all sufficiently large }n
 \quad\Longrightarrow\quad w_i(n)=0\quad\text{for every }i,n.
\]
\end{corollary}
\begin{proof}
Choose a common period $Q$. On each progression $n=a+Qt$, the coefficients
are constant and the cross determinants become $Q(a_ib_j-a_jb_i)\ne0$.
Shift beyond the threshold and apply the affine independence just proved.
\end{proof}
The same argument preserves the basis after deletion of a finite prefix.
For the two zero channels, large $n\equiv1\pmod k$ and large powers of $k$
replace the two evaluations used above.
